% Construction 1 on a multicolored clique with k = 4 colors, drawn to scale.
% Adapted from the slide "Reduction from k-Multicolored Clique".
\begin{tikzpicture}[xscale=0.42,yscale=0.8]
  \colorlet{c0}{red!85!black}
  \colorlet{c1}{blue!70!white}
  \colorlet{c2}{yellow!85!orange}
  \colorlet{c3}{green!60!black}
  \def\hh{0.19}
  \def\g{0.06}
  % \jobbar[outline options]{x1}{x2}{y}{left color}{right color}
  \newcommand{\jobbar}[6][]{%
    \shade[left color=#5,right color=#6] (#2,#4-\hh) rectangle (#3,#4+\hh);
    \draw[line width=0.35pt,#1] (#2,#4-\hh) rectangle (#3,#4+\hh);}

  % the graph
  \node[circle,draw,fill=c0,inner sep=3.2pt] (v0) at (-8.6,3.6) {};
  \node[circle,draw,fill=c1,inner sep=3.2pt] (v1) at (-7.0,6.2) {};
  \node[circle,draw,fill=c2,inner sep=3.2pt] (v2) at (-4.6,4.6) {};
  \node[circle,draw,fill=c3,inner sep=3.2pt] (v3) at (-6.2,2.0) {};
  \foreach \s/\t in {v0/v1,v0/v2,v0/v3,v1/v2,v1/v3,v2/v3} \draw[line width=0.8pt] (\s) -- (\t);
  \node at (-6.6,7.4) {\small $G$};

  % time axis and the four vertex windows
  \foreach \x in {1,7,13,19,25}
    \draw[line width=0.4pt,dotted,gray] (\x,0) -- (\x,7.6);
  \draw[line width=0.8pt,->] (-0.4,0) -- (27.2,0);
  \foreach \x in {1,7,13,19,25,26}{
    \draw[line width=0.6pt] (\x,0.12) -- (\x,-0.12);
    \node at (\x,-0.5) {\scriptsize \x};}

  % validation machine: the vertex job of each vertex for its own color
  \draw[fill=gray!30,line width=0.35pt] (-2.3,1-0.3) rectangle (-0.9,1+0.3);
  \foreach \c in {0,...,3}{
    \pgfmathsetmacro{\xa}{6*\c+1+\g}
    \pgfmathsetmacro{\xb}{6*\c+7-\g}
    \jobbar{\xa}{\xb}{1}{c\c}{c\c}}

  % one edge selection machine per color combination
  \foreach \a/\b/\row in {0/1/1,0/2/2,0/3/3,1/2/4,1/3/5,2/3/6}{
    \pgfmathsetmacro{\y}{1+\row}
    \pgfmathsetmacro{\su}{6*(\a+1)-\b-1}
    \pgfmathsetmacro{\sv}{6*(\b+1)-\a-1}
    \shade[left color=c\a,right color=c\b] (-2.3,\y-0.3) rectangle (-0.9,\y+0.3);
    \draw[line width=0.35pt] (-2.3,\y-0.3) rectangle (-0.9,\y+0.3);
    \begin{scope}[opacity=0.45]
      \jobbar{1+\g}{\su-\g}{\y}{c\b}{c\a}
      \jobbar{\sv+1+\g}{26-\g}{\y}{c\b}{c\a}
    \end{scope}
    \jobbar{\su+\g}{\su+1-\g}{\y}{c\a}{c\b}
    \jobbar{\sv+\g}{\sv+1-\g}{\y}{c\b}{c\a}
    \jobbar[line width=1.1pt]{\su+1+\g}{\sv-\g}{\y}{c\a}{c\b}}
\end{tikzpicture}
